% !TEX root = ../main.tex
\section{Introduction}

Linear programming has been the workhorse of strategic forest planning for half a century. Since the Model I and Model II formulations of \citet{johnson1977techniques}, LP harvest-scheduling models---most prominently the USDA Forest Service's FORPLAN system and its successors---have been used to project long-term timber supply, to test harvest-flow policies, and to support trade-off analysis through the shadow prices that an optimal LP solution provides. In nearly all of this practice, the model is solved once, over the full planning horizon, as an \emph{open-loop} problem: the entire sequence of future decisions is chosen at time zero, and the resulting schedule of harvest levels and residual forest structure is presented as the plan.

There is a trap in this practice, and it is a deep one. An open-loop LP solution is a plan about the future made as if the planner could precommit to every future decision. But real planning is sequential: each period, the planner observes the realized forest state and re-solves the planning problem from that state, under the same objectives and constraints. If the tail of the original plan is not optimal for the subproblem that starts at the realized state, the re-solving planner will deviate from it. The announced plan then fails Bellman's principle of optimality (\citep{bellman1957dynamic}): it is \emph{dynamically inconsistent}, and the harvest trajectory, residual growing stock, and shadow prices it projects are not credible statements about what will or should happen. A plan that will not be followed cannot be a sound basis for policy, and comparative analysis performed on it---including the shadow-price trade-off analysis that motivates much of the modelling---is biased in ways that are difficult to bound.

This trap was demonstrated carefully and thoroughly more than three decades ago. In his 1991 PhD thesis at the University of California, Berkeley, P. J. Daugherty proved analytically and demonstrated numerically that LP-based forest-planning models solved as open-loop formulations admit dynamically inconsistent plans, identified the structural conditions under which the inconsistency arises, and characterized its behaviour over a wide range of initial forest conditions and harvest policies \citep{daugherty1991credibility}. The thesis remains, to our knowledge, the most complete treatment of the problem in the forest-planning literature. It is also effectively inaccessible: it was never published in the peer-reviewed literature, it exists only as an archival hard copy, and it can be obtained only as a custom print-on-demand from the UC Berkeley Library at a cost of roughly \$200. The phenomenon itself has surfaced in the literature under various names, before and since---\citet{mcquillan1986declining} named the ``declining even-flow effect'' in national forest planning; \citet[pp.~331--332]{gunn2007models}, citing \citet{daugherty1991credibility}, calls it the ``declining non-declining yield''; and \citet{paradis2013risk} documented the systematic drift that arises between planned and implemented harvest under incoherent hierarchical planning, citing \citet{daugherty1991credibility} directly and using ``credibility'' in the thesis's sense. Yet the original analysis stays unread: the field keeps rediscovering pieces of the phenomenon, typically without access to, or citation of, the original analysis.

This paper exists to fix the accessibility problem, not to claim new science. We present an open, transparent, and fully reproducible reproduction of the modelling methods and case study of \citet{daugherty1991credibility}, built on the open-source \texttt{ws3} wood-supply model \citep{paradis2026ws3}, together with this paper. Every element of the reproduction---the case-study data extracted from the Umpqua planning documents, the LP formulations, the sequential-replanning simulator, the experiment grid, and the validation record---is available in a public repository with documented provenance and a test suite. We emphasize the framing: the science reported here is a faithful reproduction of \citet{daugherty1991credibility}. Our only novel content is the open, reproducible software stack and a citable reference. The case-study data are extracted from the public-domain Umpqua planning documents (the 1987 Draft EIS and the 1990 FEIS Appendix B yield tables and economics).

The remainder of this paper is organized as follows. Section 2 reviews the theory of dynamic inconsistency and the open-loop LP forest-planning problem, including the Model I/II/III formulation distinction. Section 3 describes the case study, the case-study data, the sequential-replanning simulator, and the experiment design. Section 4 presents the reproduced results. Section 5 discusses the drivers of the inconsistency, a subtle and counter-intuitive nuance concerning the discount rate, implications for the literature, remedies, and limitations. Section 6 concludes.
